\documentclass[12pt,a4paper]{article}
\usepackage{luatexja}
\usepackage{luatexja-fontspec}
\usepackage{fontspec}
\usepackage{graphicx}
\usepackage{listings}
\usepackage[margin=2.4cm]{geometry}
\usepackage[hidelinks]{hyperref}
\usepackage{amsmath,amssymb}
\usepackage{booktabs}
\usepackage{array}
\usepackage{xcolor}
\usepackage{enumitem}
\usepackage{float}
\usepackage{caption}

\lstset{
    basicstyle=\ttfamily\small,
    breaklines=true,
    frame=single,
    backgroundcolor=\color{gray!6},
    columns=fullflexible,
}

\graphicspath{{images/}}
\setmainjfont{Noto Serif CJK SC}[
    BoldFont=Noto Serif CJK SC Bold,
    ItalicFont=Noto Serif CJK SC Light,
]
\setmainfont{Times New Roman}

\linespread{1.3}
\setlength{\parindent}{2em}
\captionsetup{font=small,labelfont=bf}
\setlist{nosep,leftmargin=2em}

\newcommand{\screenshot}[2]{%
  \IfFileExists{images/#1}%
    {\includegraphics[width=0.92\linewidth]{#1}}%
    {\fbox{\parbox[c][0.26\textheight][c]{0.86\linewidth}{\centering
      \textbf{待补截图}\\[4pt] #2}}}%
}

\title{\textbf{深度伪造检测系统\\使用手册}}
\author{课程设计小组}
\date{\today}

\begin{document}
\maketitle

\begin{abstract}
\noindent
本手册面向系统的使用者，介绍深度伪造检测系统的安装、启动、Web 与 API 使用、
批量检测、配置项与常见问题。系统输入一张人脸图片，输出伪造概率、Grad-CAM 热图、
频谱图，并可生成带哈希链审计的 PDF 报告。
\end{abstract}

\tableofcontents
\newpage

\section{系统概述}

\subsection{功能}

\begin{table}[H]
\centering
\caption{系统功能}
\small
\begin{tabular}{p{4.2cm}p{10cm}}
\toprule
功能 & 说明 \\
\midrule
图片检测 & 上传 JPG/PNG/BMP/WebP，返回伪造概率与判定 \\
人脸裁剪 & 自动检测最大人脸并外扩裁剪；检测不到时使用整图 \\
可解释 & Grad-CAM 热图 + log$|$FFT$|$ 频谱图 \\
PDF 报告 & 单页 A4，含判定、元信息、三张图、审计哈希 \\
审计日志 & JSONL 哈希链，可用接口校验完整性 \\
API & REST 接口，支持脚本调用 \\
批量检测 & 命令行扫描目录，输出 CSV \\
\bottomrule
\end{tabular}
\end{table}

\subsection{模型与指标}

主模型为 ResNet18（ImageNet 预训练）RGB 基线，训练使用 FaceForensics++ C23 的 580 个身份。
在 3,600 帧测试集上：帧级 AUC 0.933、TPR@FPR=1\% 0.596；视频级 AUC 0.950、TPR@1\% 0.700；
跨库 Celeb-DF AUC 0.716。

系统内置两个阈值：\textbf{严格模式}（验证集 FPR$\le$1\%，约 0.9974）与\textbf{演示模式}
（验证集 F1 最优，约 0.0962）。阈值随模型权重保存在 checkpoint 中。

\subsection{系统结构}

\begin{itemize}
  \item \texttt{app/detector.py}：模型加载、人脸裁剪、打分、Grad-CAM、频谱；
  \item \texttt{app/main.py}：FastAPI 路由；
  \item \texttt{app/report.py}：Jinja2 + WeasyPrint 生成 PDF；
  \item \texttt{app/audit.py}：审计哈希链；
  \item \texttt{app/templates/}：Web 页面与报告模板。
\end{itemize}

\section{运行环境}

\subsection{硬件}

\begin{itemize}
  \item CPU：4 核以上；本机实测 \textbf{AMD Ryzen 9 7945HX with Radeon Graphics}（16 核 32 线程）；
  \item 内存：8 GB 以上；
  \item 磁盘：5 GB 以上（含模型权重与依赖）；
  \item GPU：可选。NVIDIA 显卡 + CUDA 驱动可显著加速（本机实测 RTX 4060 Laptop）。
\end{itemize}

\subsection{软件}

\begin{table}[H]
\centering
\caption{软件依赖}
\small
\begin{tabular}{ll}
\toprule
项 & 要求 \\
\midrule
操作系统 & Linux（Arch Linux 实测）\\
Python & 3.12 及以上（实测 3.14）\\
PyTorch & 2.15（GPU 版 cu134 或 CPU 版均可）\\
Web & FastAPI + Uvicorn \\
PDF & WeasyPrint（依赖 pango/cairo）\\
字体 & Noto Serif/Sans CJK（PDF 中文）\\
\bottomrule
\end{tabular}
\end{table}

\section{安装}

\subsection{创建虚拟环境}

\begin{lstlisting}
python3 -m venv .venv
source .venv/bin/activate
\end{lstlisting}

\subsection{安装依赖}

GPU（cu134 轮子）：

\begin{lstlisting}
uv pip install torch torchvision \
  --index-url https://download.pytorch.org/whl/cu134
\end{lstlisting}

CPU：

\begin{lstlisting}
uv pip install torch torchvision \
  --index-url https://download.pytorch.org/whl/cpu
\end{lstlisting}

Web 与 PDF 依赖：

\begin{lstlisting}
uv pip install fastapi "uvicorn[standard]" jinja2 python-multipart \
  pillow weasyprint numpy opencv-python-headless scikit-learn pyyaml tqdm
\end{lstlisting}

\subsection{准备模型文件}

需要两个文件：

\begin{lstlisting}
runs/rgb_700/best.pt                          # 主模型权重（约 44 MB）
models/face_detection_yunet_2023mar.onnx      # 人脸检测模型（约 228 KB）
\end{lstlisting}

YuNet 缺失时程序会尝试自动下载；离线环境请提前放置。

\subsection{自检}

\begin{lstlisting}
python -c "import torch; print(torch.__version__, torch.cuda.is\_available())"
\end{lstlisting}

输出第二项为 \texttt{True} 表示 GPU 可用。

\section{启动与停止}

\subsection{本机访问}

\begin{lstlisting}
python -m uvicorn app.main:app --host 127.0.0.1 --port 8000
\end{lstlisting}

浏览器打开 \url{http://127.0.0.1:8000}。

\subsection{局域网访问}

\begin{lstlisting}
python -m uvicorn app.main:app --host 0.0.0.0 --port 8000 --workers 1
\end{lstlisting}

其他机器访问 \url{http://<本机IP>:8000}，需放行防火墙 8000 端口。

\textbf{注意}：\texttt{--workers} 保持为 1。模型常驻显存，多 worker 会复制多份模型。

\subsection{后台常驻}

开发调试可直接前台运行；长期运行建议用 systemd（见 \texttt{deploy/README.md}）或：

\begin{lstlisting}
nohup python -m uvicorn app.main:app --host 0.0.0.0 --port 8000 \
  > data/app_server.log 2>&1 &
\end{lstlisting}

\subsection{停止}

前台运行按 \texttt{Ctrl+C}；systemd 方式：

\begin{lstlisting}
sudo systemctl stop deepfake
\end{lstlisting}

\subsection{环境变量}

\begin{table}[H]
\centering
\caption{环境变量}
\small
\begin{tabular}{lll}
\toprule
变量 & 默认值 & 说明 \\
\midrule
\texttt{APP\_CKPT} & \texttt{runs/rgb\_700/best.pt} & 模型权重路径 \\
\bottomrule
\end{tabular}
\end{table}

\section{Web 界面使用}

\subsection{操作步骤}

\begin{enumerate}
  \item 打开首页，拖拽图片到虚线框，或点击“选择图片”；
  \item 点击“开始检测”；
  \item 查看伪造概率、判定结果、Grad-CAM 热图与频谱图；
  \item 点击“下载 PDF 审计报告”保存报告。
\end{enumerate}

\begin{figure}[H]
\centering
\screenshot{fig_ui.png}{Web 界面：上传区、分数、判定、热图、频谱、PDF 下载}
\caption{Web 界面}
\label{fig:ui}
\end{figure}

\subsection{结果解读}

\begin{table}[H]
\centering
\caption{结果字段}
\small
\begin{tabular}{p{3.2cm}p{10.5cm}}
\toprule
字段 & 含义 \\
\midrule
伪造概率 & 模型输出的 $p(\text{fake})$，0--1 \\
判定 & 概率高于当前阈值判“伪造”，否则“真实” \\
使用阈值 & 验证集 F1 最优的固定阈值（约 0.0962） \\
严格阈值 & FPR$\le$1\% 的阈值（约 0.9974），仅用于对照 \\
Grad-CAM & 模型关注区域，越红表示贡献越大 \\
频谱图 & log$|$FFT$|$，伪彩色；亮点表示该频率成分强 \\
审计哈希 & 本条记录在哈希链中的 SHA-256 \\
\bottomrule
\end{tabular}
\end{table}

\subsection{阈值说明}

系统使用模型权重中固定保存的阈值（验证集 F1 最优，约 0.0962），Web 界面不提供阈值调节。
模型权重中还保存了满足 FPR$\le$1\% 的严格阈值（约 0.9974），仅用于结果对照。
两个阈值均由验证集确定；跨数据集时分数分布可能偏移，阈值的适用性会下降。

\subsection{PDF 报告}

点击下载后生成单页 A4 报告，包含：报告编号、时间、模型、判定与概率、文件 SHA-256、
两个阈值、人脸框、输入尺寸、推理耗时、审计哈希，以及原图、热图、频谱三张图。

\begin{figure}[H]
\centering
\screenshot{fig_pdf.png}{PDF 报告第一页}
\caption{PDF 报告}
\label{fig:pdf}
\end{figure}

\section{API 使用}

\subsection{接口列表}

\begin{table}[H]
\centering
\caption{HTTP 接口}
\small
\begin{tabular}{lll}
\toprule
方法 & 路径 & 说明 \\
\midrule
GET & \texttt{/} & Web 页面 \\
POST & \texttt{/v1/detect/image} & 图片检测 \\
GET & \texttt{/v1/report/\{id\}.pdf} & 下载 PDF \\
GET & \texttt{/v1/audit/verify} & 校验审计链 \\
GET & \texttt{/healthz} & 健康检查 \\
\bottomrule
\end{tabular}
\end{table}

\subsection{调用示例}

检测一张图片：

\begin{lstlisting}
curl -X POST "http://127.0.0.1:8000/v1/detect/image" \
     -F "file=@test.jpg"
\end{lstlisting}

返回（节选）：

\begin{lstlisting}
{
  "id": "a582ee564c1c",
  "score": 0.999999,
  "verdict": "fake",
  "threshold": 0.0913,
  "threshold_strict": 0.9996,
  "latency_ms": 10.9,
  "artifacts": {
    "heatmap": "/artifacts/a582ee564c1c/heatmap.png",
    "spectrum": "/artifacts/a582ee564c1c/spectrum.png",
    "pdf": "/v1/report/a582ee564c1c.pdf"
  },
  "audit_hash": "cfca33..."
}
\end{lstlisting}

下载报告、校验审计：

\begin{lstlisting}
curl -o report.pdf http://127.0.0.1:8000/v1/report/a582ee564c1c.pdf
curl http://127.0.0.1:8000/v1/audit/verify
# {"ok": true, "n": 2}
\end{lstlisting}

\subsection{参数}

\begin{table}[H]
\centering
\caption{\texttt{/v1/detect/image} 参数}
\small
\begin{tabular}{lll}
\toprule
参数 & 位置 & 说明 \\
\midrule
\texttt{file} & form-data & 图片文件，必填 \\
\bottomrule
\end{tabular}
\end{table}

\subsection{返回字段}

\begin{table}[H]
\centering
\caption{返回字段}
\small
\begin{tabular}{p{3.4cm}p{10.4cm}}
\toprule
字段 & 含义 \\
\midrule
\texttt{id} & 本次检测编号 \\
\texttt{score} & 伪造概率 \\
\texttt{verdict} & \texttt{fake} / \texttt{real} \\
\texttt{threshold} & 本次使用的阈值 \\
\texttt{threshold\_strict} & 严格阈值（FPR$\le$1\%） \\
\texttt{face\_box} & 人脸框 $[x,y,w,h]$，未检测到为 null \\
\texttt{latency\_ms} & 本次推理耗时 \\
\texttt{artifacts} & 热图、频谱、PDF 的 URL \\
\texttt{audit\_hash} & 审计记录哈希 \\
\bottomrule
\end{tabular}
\end{table}

\subsection{错误码}

\begin{table}[H]
\centering
\caption{常见错误}
\small
\begin{tabular}{lll}
\toprule
状态码 & 原因 & 处理 \\
\midrule
400 & 空文件 / 无法解码 & 检查图片格式 \\
404 & 报告 ID 不存在 & 确认 ID \\
422 & 缺少 file 字段 & 检查 form-data \\
500 & 服务内部错误 & 查看日志 \\
\bottomrule
\end{tabular}
\end{table}

\section{批量检测}

对目录下所有图片打分，不启动服务：

\begin{lstlisting}
python scripts/07_batch_detect.py --input data/frames/Real --limit 100
\end{lstlisting}

输出 CSV（\texttt{docs/experiments/batch.csv}）：

\begin{lstlisting}
path,score,verdict,latency_ms
data/frames/Real/00015.jpg,0.000019,real,36.8
...
\end{lstlisting}

参数：\texttt{--input} 目录，\texttt{--limit} 限制数量，\texttt{--ckpt} 指定权重，
\texttt{--mode f1|strict} 选择阈值。

\section{数据与配置}

\subsection{目录结构}

\begin{lstlisting}
data/app/
├── uploads/    上传原图
├── reports/    热图、频谱、PDF
├── results/    每次检测的 JSON
└── audit.jsonl 审计日志（哈希链）
\end{lstlisting}

\subsection{审计日志}

每行一条 JSON，含 \texttt{id}、\texttt{sha256}、\texttt{score}、\texttt{verdict}、
\texttt{prev\_hash}、\texttt{hash}。篡改任意一条，后续哈希校验失败。

\subsection{清理}

上传图、热图、PDF 可定期清理；审计日志建议保留。清空后哈希链从头开始。

\subsection{更换模型}

把新权重放到 \texttt{runs/} 下，启动时设置 \texttt{APP\_CKPT} 指向它即可；阈值随 checkpoint 读取。

\section{常见问题}

\begin{enumerate}
  \item \textbf{首页 500}：Starlette 1.x 的 \texttt{TemplateResponse} 需要 \texttt{request} 作为第一个参数，本项目已修复；旧版本请升级代码。
  \item \textbf{端口被占用}：换 \texttt{--port 8001}，或 \texttt{lsof -i:8000} 查占用进程。
  \item \textbf{CUDA 不可用}：检查驱动与 \texttt{torch.cuda.is\_available()}；CPU 也能跑（约 37 ms/张）。
  \item \textbf{PDF 中文乱码}：安装 Noto CJK 字体（\texttt{fonts-noto-cjk}）。
  \item \textbf{检测不到人脸}：程序会用整图；正面、清晰的人脸效果最好。
  \item \textbf{上传失败}：检查大小与格式；Nginx 需调大 \texttt{client\_max\_body\_size}。
  \item \textbf{磁盘增长快}：定期清理 \texttt{data/app/} 下的旧文件。
\end{enumerate}

\section{安全与合规}

\begin{itemize}
  \item 建议本地/内网部署，避免上传敏感人脸到公网；
  \item 审计日志包含文件哈希，注意隐私与留存期限；
  \item 检测结果仅供参考，不构成事实认定；高利害场景应人工复核；
  \item 使用需遵守《互联网信息服务深度合成管理规定》等法规。
\end{itemize}

\section{性能}

\begin{table}[H]
\centering
\caption{单图耗时（含人脸检测与 Grad-CAM）}
\small
\begin{tabular}{lcc}
\toprule
设备 & 中位数 & p95 \\
\midrule
RTX 4060 Laptop & 10.9 ms & -- \\
CPU（Ryzen 9 7945HX） & 36.8 ms & 42.0 ms \\
\bottomrule
\end{tabular}
\end{table}

PDF 生成约 0.2--0.5 s，单页约 500 KB。

\appendix
\section{目录结构}

\begin{lstlisting}
deepfake-det/
├── app/            服务代码
├── src/            模型、数据、训练、评测
├── configs/        训练配置
├── runs/           模型权重与评测结果
├── scripts/        数据处理、批量、鲁棒性、可解释性脚本
├── deploy/         Dockerfile、systemd、nginx
├── data/           数据集与运行时数据
└── report/         LaTeX 报告与使用手册
\end{lstlisting}

\section{响应示例}

完整返回：

\begin{lstlisting}
{
  "id": "a582ee564c1c",
  "ts": "2026-11-20T23:17:46",
  "filename": "00015.jpg",
  "sha256": "c02a7ec3...",
  "score": 0.999999,
  "threshold": 0.0913,
  "threshold_strict": 0.9996,
  "verdict": "fake",
  "label": 1,
  "face_box": [4, 0, 250, 250],
  "latency_ms": 10.9,
  "input_size": [256, 256],
  "ckpt": "runs/rgb_700/best.pt",
  "artifacts": {
    "input": "/artifacts/a582ee564c1c/input.jpg",
    "heatmap": "/artifacts/a582ee564c1c/heatmap.png",
    "spectrum": "/artifacts/a582ee564c1c/spectrum.png",
    "pdf": "/v1/report/a582ee564c1c.pdf"
  },
  "audit_hash": "cfca3363..."
}
\end{lstlisting}

\end{document}
